\documentclass{article}
\usepackage{amsmath, amsthm, amsfonts, amssymb}
\newcommand{\op}{\prime}
\DeclareMathOperator{\re}{Re}
\DeclareMathOperator{\im}{Im}
\DeclareMathOperator{\curl}{curl}
\DeclareMathOperator{\res}{Res}
\begin{document}
\title{Some Contour Integrals}
\author{Amey Joshi}
\maketitle
\begin{enumerate}
\item Let
\begin{equation}\label{e1}
I = \int_0^\infty \frac{dx}{(1 + x^n)} \text{ where } n \in \mathbb{N}, n \ge 2.
\end{equation}
The zeroes of $z^n + 1 = 0$ are at $z = \exp(i\pi/n + 2\pi m i/n)$ for $m = 0,
1, \ldots, n - 1$. They are separated by an angle of $2\pi/n$. Consider a wedge
contour formed by considering the sector of a circle of angle $2\pi/n$ and 
radius $R$ whose one arm along the real axis. It encloses only one pole $z = 
\exp(i\pi/n)$. The integral along the contours arc vanishes as $R \rightarrow
\infty$. If $S_R$ denotes the wedge contour then,
\begin{eqnarray*}
\oint_{S_R}\frac{dz}{1 + z^n} &=& \int_0^R\frac{dx}{1 + x^n} + 
\int_R^0\frac{e^{2\pi i/n}dr}{1 + r^n} \\
&=& (1 - e^{2\pi i/n})\int_0^R\frac{dx}{1 + x^n}
\end{eqnarray*}
If $f(z) = 1/(1 + z^n)$ then the left hand side can be evaluated by residue
theorem to be
\[
2\pi i\res(f; e^{\pi i/n}) = 2\pi i\frac{1}{nz^{n-1}}\Big|_{e^{\pi i/n}}
= \frac{2\pi i}{ne^{(n-1)\pi i/n}} = \frac{-2\pi i}{n}e^{i\pi/n}.
\]
Therefore,
\[
I = -\frac{2\pi i}{n}\frac{e^{i\pi/n}}{1 - e^{2\pi i/n}} = 
-\frac{\pi}{n}\frac{2\pi}{e^{-i\pi/n} - e^{\pi i/n}}
\]
or that
\begin{equation}\label{e2}
\int_0^\infty \frac{dx}{1 + x^n} = \frac{(\pi/n)}{\sin(\pi/n)}.
\end{equation}
\end{enumerate}
\end{document}
